\manualchapter{Overview and first patch}{sec:quick-start}
CV Genie writes control voltages into game memory in RackNES. Eight input
rows each target one named element from a built-in game map. Depending on
the element, voltage continuously sets a byte value or a trigger toggles a
stored 0/1 state.

In the module browser, choose \textbf{CV Genie (Input)}. This version
provides the input expander only. It does not output game-memory voltages,
accept typed Game Genie codes, or offer an arbitrary-address editor.

\subsection{Set up in this order}
\begin{enumerate}
\item Load and start a supported game in RackNES. CV Genie includes maps for
\textbf{Super Mario Bros.} and \textbf{The Legend of Zelda}.
\item Place CV Genie \textbf{immediately to the right of RackNES}, touching
it in the same row. There must be no module or empty space between them.
\item Leave the eight inputs unplugged. Click the top display and select
the game that is actually running in RackNES.
\item Click a row's display and choose a \textbf{Game Element}. For a first
Super Mario Bros. experiment, try \textbf{Player Horizontal Screen Position}.
\item Connect an adjustable \textbf{0--10 V} source to that row. Start near
0 V and move it slowly while watching the game. This writes a screen-position
byte; game logic decides how it appears on screen.
\item Disconnect that input to stop writing. Other rows can be assigned and
connected in the same way, one at a time.
\end{enumerate}

\subsection{An important setup limitation}
\textbf{Choose the game before opening row menus, and assign a row before
connecting CV.} The current implementation does not guard every access to
an unselected game or an Unassigned row. These combinations can access
invalid memory and may crash Rack. An Unassigned row with a cable is not a
safe disabled input; unplug it first.

\subsection{What a game map means}
The top selector chooses a table of addresses. It neither loads a ROM nor
detects which game RackNES is running. Map labels describe intended game
variables, but the game can overwrite those variables or interpret them
differently in another scene or ROM revision. Make a RackNES SAVE snapshot
before experimenting and disconnect Genie inputs before restoring it.
